\section{Runge--Kutta 与自适应步长}
\label{sec:08-Numerical-Analysis/06-ODE/02}

拿着上一节 Euler 的经验回到 \(y'=-y,\;y(0)=1\)。Euler 用 \(h=0.5\) 在 \(t=1\) 时误差约 0.118，用 \(h=0.1\) 才能把误差压到可接受的范围。问题是：减小步长意味着多走十倍步数，每一步都在做浮点运算、积累舍入误差。能不能不缩小步长，而是在一步之内多获取一些信息，走得更准？

\subsection{看两次斜率，胜过看一次}

Euler 只看起点的斜率，然后沿着它直走到终点。一个自然的改进是：先用起点斜率走到中途，在中途重新评估一次斜率，再用这个\"路况更新后的方向\"走完全程。这就是显式中点法（也叫 RK2）：

\[
k_1=f(t_n,y_n),
\]
\[
k_2=f\left(t_n+\frac h2,\;
y_n+\frac h2k_1\right),
\]
\[
y_{n+1}=y_n+hk_2.
\]

为什么这样做就比 Euler 高一阶？把精确解在 \(t_n+h/2\) 处展开：
\[
y(t_n+h)=y(t_n)+h\,y'\!\left(t_n+\tfrac h2\right)+O(h^3).
\]
中点斜率乘整步长，局部误差已经是 \(O(h^3)\) 而非 Euler 的 \(O(h^2)\)。但我们并没有直接算真实中点的斜率——我们用 Euler 半步的结果预测了一个近似中点。这个预测误差是 \(O(h^2)\)，但经过 Lipschitz 连续的 \(f\) 传递、再被外层因子 \(h\) 乘上之后，局部误差仍然是 \(O(h^3)\)。于是局部三阶累积成全局二阶——和上一节\"局部高一阶、全局低一阶\"是一模一样的机制。

Runge--Kutta（RK）方法把这个思路系统化了：一步之内，在若干个位置评估斜率，然后加权组合，用不同位置的斜率信息互相校正，抵消掉低阶误差项。

\subsection{一般 RK 与 Butcher 记号}

一个 \(s\) 阶段的 RK 方法写成：
\[
k_i=f\left(t_n+c_ih,\;
y_n+h\sum_{j=1}^sa_{ij}k_j\right),\qquad i=1,\dots,s,
\]
\[
y_{n+1}=y_n+h\sum_{i=1}^sb_ik_i.
\]

所有的系数——\"几点去评估\"（\(c_i\)）、\"评估结果怎么互相影响\"（\(a_{ij}\)）、\"最后怎么加权\"（\(b_i\)）——都装进一个叫 Butcher tableau 的表里。显式 RK 满足 \(a_{ij}=0\) 对 \(j\ge i\)，各阶段按顺序直接计算，不需要解方程；隐式 RK 的阶段之间互相耦合，每步要解非线性系统。

\subsection{经典 RK4 为什么能到四阶}

经典的四阶 RK 在一个时间步内评估四次斜率：起点一次，用起点斜率预测的两个中点各一次，再用中点斜率预测的终点一次。四次评估结果分别以权重
\[
\frac16(1,2,2,1)
\]
组合——中间两个中点的权重各是两端的二倍，反映了中点信息在整个区间中承载更高的代表权重。

它的局部截断误差是 \(O(h^5)\)，全局是 \(O(h^4)\)。四个阶段不自动意味着四阶——阶数来自系数精确满足一组 Taylor 展开匹配条件（或等价地，Butcher 的 rooted-tree 阶条件）。随便取四个斜率再随便加权，几乎得不到四阶精度，甚至可能比 Euler 的一阶还差。

高阶方法减少达到固定精度所需的步数，但每一步的函数评估次数也更多，稳定区域也不一样。比较两个方法时，\"阶数\"只是众多指标之一：总函数评估次数、是否容易实现自适应步长、对刚性问题的稳定区域、内存占用、事件处理复杂度——都要放进天平。

\subsection{固定步长不知道轨迹哪里变陡}

一条 ODE 的轨迹可能长时间平缓，只在很短的区间内剧烈变化。固定步长面临一个两难：用小步长走完全程，在平缓段白白浪费评估次数；用大步长走完全程，在陡峭段又会跨过去而完全失准。

嵌入式 RK 方法用了一个巧妙的办法：同一组阶段评估结果，用两套不同的权重 \(b_i\) 和 \(\hat b_i\) 分别组合，得到两个近似——一个是 \(p\) 阶，另一个是 \(p-1\) 阶（或 \(p+1\) 阶）。两个近似的差就是当前步局部截断误差的估计，不需要额外评估函数。

标准做法是先算出每个分量的误差，再用缩放误差统一判断这一整步能不能被接受。对含 \(d\) 个分量的状态向量：
\[
E=\left[
\frac1d\sum_{j=1}^d
\left(
\frac{e_j}{a_{\rm tol}+r_{\rm tol}\max(|y_{n,j}|,|y_{n+1,j}|)}
\right)^2
\right]^{1/2}.
\]
分母里同时放了绝对容差 \(a_{\rm tol}\) 和相对容差 \(r_{\rm tol}\)。\(a_{\rm tol}\) 负责\"当分量接近零时，绝对大小重要\"——你不能因为一个量本身是 \(10^{-13}\) 就说它的小变化不重要。\(r_{\rm tol}\) 负责\"当分量的绝对值很大时，重要的是相对变化\"——\(y_j=10^6\) 时，误差 \(10^{-3}\) 可能只是舍入级别的。

若 \(E\le1\)，接受这一步；若 \(E>1\)，拒绝它，缩小步长重走。新步长的调整公式大致是
\[
h_{\rm new}
=\text{safety}\times h\cdot E^{-1/(p+1)},
\]
其中指数 \(1/(p+1)\) 来自局部截断误差对步长的关系，safety 是一个略小于 1 的安全系数（典型值 0.8--0.9），用于防止刚刚勉强通过的一步引发下一轮拒绝。步长的放大和缩小通常还设上限和下限，避免步长在各步之间剧烈震荡。

\subsection{局部控制住了，不等于全局就对了}

嵌入式估计只反映当前这一步、当前这个解附近的局部截断误差。它背后有两个重要前提：第一，方法处于渐近区（\(h\) 已经足够小，使得截断阶的渐近行为成立）；第二，被控制的只是时间离散误差。

非线性动力学会把局部误差放大或衰减——这和 Euler 误差分析中的 Grönwall 放大是同一件事。函数不光滑、事件附近、遇到刚性问题时，实际阶数可能达不到理论阶数——这种现象叫 order reduction。此时嵌入式误差估计的量级关系不再成立，估计值可能严重失真。

所以局部容差 \(a_{\rm tol}, r_{\rm tol}\) 和最终解的全周期精度之间只是相关关系，不是等号。把 tol 设成 \(10^{-6}\)，不能理解为\"最终答案误差必小于 \(10^{-6}\)\"。初始值不准确、模型参数有误差、物理模型本身做了近似——这些都不会因为你把步长容差调小就自动消失。

\subsection{事件不是离散采样能刚好踩到的}

碰到\"物体何时撞到地面\"（\(g(t,y)=0\)）或\"开关何时触发\"这类事件时，按步长离散走大概率会刚好跨过事件点。成熟求解器在一个被接受的步内，用已经评估过的斜率构造一个稠密输出插值多项式，再以这个多项式作为连续近似，对事件函数求根，找到更精确的事件时间。跨过事件后可能需要切换动力系统（比如碰撞反弹后的新方程），此时必须用新模型的初值重启积分——拿着跨过不连续面之前的旧信息继续用高阶多项式外推，等于假装不连续不存在。

\subsection{RK 不是万能的}

经典显式 RK 在光滑、非刚性、步长合理的初值问题中非常可靠。但以下场景需要退一步选择更合适的方法：

\begin{itemize}
\tightlist
\item
  强刚性：显式 RK 的稳定区域有限，步长会被稳定性而非精度压到极小。BDF、Radau 等隐式方法虽然每步更贵，但允许大得多的步长——\hyperref[sec:08-Numerical-Analysis/06-ODE/03]{``稳定区域与刚性方程''}会给出定量判据。
\item
  Hamilton 系统长时间积分：能量和辛结构的保持比局部截断阶更重要。辛 RK 方法或 Verlet 类方法专门为此设计——它们在十万步之后仍然如实反映系统的定性行为，而经典 RK4 可能已经漂出了相空间的真实区域。
\item
  强不连续或随机驱动：需要匹配对不连续敏感的专用方法，或应对随机项的随机数值格式。
\item
  极大系统：Jacobian 矩阵的构造和线性求解成本太高时，Jacobian-free、指数积分或算子分裂可能更适合。
\end{itemize}

\"RK4 很准\"是一句局部成立的判断——在它适合的土地上，它确实简洁又高效。但离开这片土地，准确度就不是阶数能担保的。
